\documentclass[aspectratio=169,10pt,english]{beamer}

\input{../../../_preambulo_beamer.tex}
\input{../../../_comandos_beamer.tex}

\selectlanguage{english}

\title{MA1001B - Statistical Modeling for Decision Making}
\subtitle{Lesson 11: Bayes' Theorem and Probability Updating}
\author{}
\institute{Tecnol\'ogico de Monterrey}
\date{}

\begin{document}

\begin{frame}[plain]
  \vspace{1.2cm}
  {\Large\bfseries MA1001B - Statistical Modeling for Decision Making\par}
  \vspace{0.7cm}
  {\large Lesson 11: Bayes' Theorem and Probability Updating\par}
  \vspace{0.7cm}
  {\color{TecRojo}\rule{\textwidth}{1.2pt}}
  \vspace{0.8cm}

  MA1001B Faculty\\[0.3cm]
  Term\\[0.3cm]
  Tecnol\'ogico de Monterrey
\end{frame}

\begin{frame}{From a Reversed Conditional to an Update}
  \begin{alertblock}{Guiding question}
    How should a prior defective rate be updated after an inspection flag,
    and what happens if that prior is much smaller than $0.10$?
  \end{alertblock}

  \vspace{0.6em}
  By the end of this lesson, you can:
  \begin{itemize}
    \item separate the prior $P(D)$ from the likelihoods $P(F\mid D)$ and $P(F\mid C)$;
    \item compute $P(D\mid F)$ with Bayes' theorem;
    \item match the posterior to the path ratio $8/17$;
    \item explain why $P(F\mid D)$ is not a posterior when the base rate is ignored.
  \end{itemize}

  \vfill
  \scriptsize All lots and inspection flags used today are synthetic. No real company data are used.
\end{frame}

\begin{frame}{Prior, Likelihood, Evidence}
  \small
  Continue the Lesson 10 tree with $P(D)=0.10$, $P(F\mid D)=0.80$, and
  $P(F\mid C)=0.10$.
  \[
    P(D\mid F)=\frac{P(F\mid D)\,P(D)}{P(F\mid D)\,P(D)+P(F\mid C)\,P(C)}.
  \]

  \vspace{0.4em}
  \begin{center}
    \begin{tabular}{lcc}
      \toprule
      \textbf{Piece} & \textbf{Symbol} & \textbf{Role} \\
      \midrule
      Prior & $P(D)=0.10$ & Rate before the flag \\
      Likelihood & $P(F\mid D)=0.80$ & Detection given defective \\
      False-flag rate & $P(F\mid C)=0.10$ & Flag given clean \\
      Evidence & $P(F)$ & Total probability of a flag \\
      \bottomrule
    \end{tabular}
  \end{center}
\end{frame}

\begin{frame}[fragile]{Step 1: Prior and Likelihoods}
  \begin{columns}[T]
    \begin{column}{0.42\textwidth}
      \small
      \begin{lstlisting}
prior = 0.10
p_f_given_d = 0.80
p_f_given_c = 0.10
      \end{lstlisting}

      \begin{block}{Verified output}
        $P(D)=0.100000$\\
        $P(C)=0.900000$\\
        $P(F\mid D)=0.800000$\\
        $P(F\mid C)=0.100000$
      \end{block}
    \end{column}
    \begin{column}{0.58\textwidth}
      \centering
      \includegraphics[width=\textwidth]{../figures/bayes_01_prior_likelihood.png}
      \vspace{0.2em}
      \scriptsize Prior versus flag likelihoods from Step 1
    \end{column}
  \end{columns}
\end{frame}

\begin{frame}{Bayes Combines the Two Flagged Paths}
  \small
  Numerator: the defective-and-flagged path
  \[
    P(D\cap F)=0.80\times 0.10=0.08.
  \]
  Denominator: law of total probability
  \[
    P(F)=0.08+0.90\times 0.10=0.17.
  \]
  Posterior:
  \[
    P(D\mid F)=\frac{0.08}{0.17}=\frac{8}{17}=0.470588.
  \]
\end{frame}

\begin{frame}[fragile]{Step 2: Posterior from Bayes' Theorem}
  \begin{columns}[T]
    \begin{column}{0.40\textwidth}
      \small
      \begin{lstlisting}
num = p_f_d * p_d
den = num + p_f_c * p_c
post = num / den
      \end{lstlisting}

      \begin{block}{Verified output}
        Numerator $=0.080000$\\
        $P(F)=0.170000$\\
        $P(D\mid F)=0.470588$\\
        Matches $8/17$
      \end{block}
    \end{column}
    \begin{column}{0.60\textwidth}
      \centering
      \includegraphics[width=\textwidth]{../figures/bayes_02_posterior.png}
      \vspace{0.2em}
      \scriptsize Numerator, $P(F)$, and posterior from Step 2
    \end{column}
  \end{columns}
\end{frame}

\begin{frame}{Step 3: A Rare Prior Collapses the Posterior}
  \begin{columns}[T]
    \begin{column}{0.42\textwidth}
      \small
      Keep $P(F\mid D)=0.80$ and $P(F\mid C)=0.10$. Change only $P(D)$ to $0.02$.

      \vspace{0.4em}
      \[
        P(F)=0.114000,
      \]
      \[
        P(D\mid F)=0.140351.
      \]

      \vspace{0.4em}
      \textbf{Limit:} treating $P(F\mid D)=0.80$ as the posterior ignores the
      base rate. Lesson 12 asks whether the counts themselves came from a
      biased frame.
    \end{column}
    \begin{column}{0.58\textwidth}
      \centering
      \includegraphics[width=\textwidth]{../figures/bayes_03_base_rate_sensitivity.png}
      \vspace{0.2em}
      \scriptsize Posterior versus prior from Step 3
    \end{column}
  \end{columns}
\end{frame}

\begin{frame}{Concept Check}
  \begin{enumerate}
    \item Which number is the prior, and which is a likelihood: $0.10$ or $0.80$?
    \item Compute $P(F)$ when $P(D)=0.10$.
    \item Why does $P(D\mid F)=8/17$ match Bayes' theorem here?
    \item If $P(D)=0.02$, why is $P(D\mid F)$ far below $0.80$?
  \end{enumerate}

  \vspace{0.6em}
  \begin{block}{Connection to Lesson 12}
    The session can continue directly into sampling bias, coverage gaps, and
    nonsampling error on a synthetic register.
  \end{block}

  \vfill
  \scriptsize Source: OpenStax, \textit{Principles of Data Science},
  Section 3.4; \textit{Introductory Business Statistics 2e}, Section 3.4.
  CC BY-NC-SA.
\end{frame}

\appendix



\end{document}
